\documentclass[11pt,letterpaper]{article}
\usepackage[margin=.79in,headheight=15pt,headsep=16pt]{geometry}
\usepackage[T1]{fontenc}
\usepackage{lmodern,microtype,amsmath,amssymb,tikz,booktabs,tabularx,enumitem,xcolor,fancyhdr,hyperref,tcolorbox,needspace}
\tcbuselibrary{breakable}
\usetikzlibrary{arrows.meta,calc}
\definecolor{olive}{HTML}{536240}\definecolor{ink}{HTML}{273529}\definecolor{teal}{HTML}{17677A}\definecolor{gold}{HTML}{A55C21}\definecolor{pale}{HTML}{F1F4EB}\definecolor{cream}{HTML}{FCF4E8}
\hypersetup{colorlinks=true,urlcolor=teal,linkcolor=olive,pdftitle={Module 8: Graph Theory and Networks}}
\color{ink}\setlength{\parindent}{0pt}\setlength{\parskip}{6pt}\setlist[enumerate]{leftmargin=*,itemsep=5pt}
\pagestyle{fancy}\fancyhf{}\fancyhead[L]{\small MATH\& 107 / MODULE 8}\fancyhead[R]{\small Graph Theory and Networks}\fancyfoot[C]{\thepage}\renewcommand{\headrulewidth}{.3pt}
\newtcolorbox{pointbox}[1]{colback=pale,colframe=olive!65,title=\textbf{#1},fonttitle=\color{olive},arc=1.6mm,boxrule=.65pt,breakable}
\newtcolorbox{trybox}[1]{colback=cream,colframe=gold!75,title=\textbf{#1},fonttitle=\color{gold},arc=1.6mm,boxrule=.65pt,breakable}
\newcommand{\sect}[1]{\Needspace{7\baselineskip}\vspace{8pt}{\Large\bfseries\color{olive}#1}\par\vspace{1pt}\textcolor{gold}{\rule{\linewidth}{.75pt}}}
\tikzset{vtx/.style={circle,draw=olive,fill=olive!15,very thick,minimum size=8mm,inner sep=1pt,font=\bfseries}, link/.style={line width=1.45pt,draw=teal!85},newlink/.style={line width=2.7pt,draw=gold,dashed},fail/.style={line width=2pt,draw=red!65,densely dashed}}
\newcommand{\network}[1]{%
\begin{tikzpicture}[scale=.82,baseline=(current bounding box.center)]
\coordinate (A) at (0,1.6);\coordinate (B) at (1.4,3);\coordinate (C) at (1.4,.2);\coordinate (D) at (3.1,1.6);\coordinate (E) at (5.3,1.6);\coordinate (F) at (6.8,3);\coordinate (G) at (6.8,.2);
\foreach \a/\b in {A/B,A/C,B/C,B/D,C/D,E/F,E/G,F/G}{\draw[link](\a)--(\b);}
\ifnum#1=2\draw[fail](D)--(E);\else\draw[link](D)--(E);\fi
\ifnum#1=1\draw[newlink](A)--(D);\fi
\ifnum#1=3\draw[newlink](D)--(F);\fi
\ifnum#1=4\draw[newlink](D)--(F);\draw[fail](D)--(E);\fi
\foreach \n in {A,B,C,D,E,F,G}{\node[vtx] at (\n){\n};}
\end{tikzpicture}}
\begin{document}
{\small\bfseries\color{teal}MATHEMATICS IN SOCIETY / MODULE 8}\par\vspace{6pt}
{\Huge\bfseries\color{olive}How connections shape a community}\par\vspace{5pt}
{\large A practical introduction to graph theory}\par\vspace{10pt}
\begin{center}\network{0}\end{center}
\begin{pointbox}{The question for this module}
A county's community organizations share alerts through established contacts. Which locations can reach one another? What happens if an important connection fails? Where should the county add one connection to make the system more reliable?
\end{pointbox}
A network diagram helps us answer questions about \emph{relationships}. The same tools can describe footpaths, food-sharing partnerships, bus connections, and communication networks. This chapter follows one county example closely, and then tests what the mathematics does and does not tell us about real decisions.

\sect{1. Turn a real situation into a network}
Each dot represents a location. We call it a \textbf{node} (or vertex). Each line represents a direct two-way communication connection. We call it an \textbf{edge}. In mathematics, a collection of nodes and edges is a \textbf{graph}. This sort of graph is different from a coordinate graph of temperature or profit.
\begin{center}\small\begin{tabular}{@{}llll@{}}\toprule A & Emergency coordinator & B & Library\\ C & Clinic & D & School\\ E & Shelter & F & North block\\ G & Community center & & \\ \bottomrule\end{tabular}\end{center}
The nine direct connections are AB, AC, BC, BD, CD, DE, EF, EG, FG. A and D may appear close on the page, but there is \emph{no} direct A--D link. The drawing is not a geographic map; its line lengths are not distances or transmission times.

The \textbf{degree} of a node counts its direct neighbors. The school D has neighbors B, C, and E, so its degree is 3. The degrees of A through G are 2, 3, 3, 3, 3, 2, 2. Their sum is 18, twice the number of edges: each edge gets counted at both ends. This is the \textbf{handshake rule}.
\begin{trybox}{Try it: community partners}
A pantry P works directly with a library L and a school S. The school also works with a community fridge F, and F works with a senior hub H. If L and S also coordinate, the edges are PL, PS, LS, SF, FH. How many direct partners does S have? Which organization has only one partner? Why might that matter?
\end{trybox}

\sect{2. Follow paths instead of guessing from a picture}
A \textbf{path} is a sequence of connections between locations. A--B--D--E--G is a path from A to G. It names five locations but travels along \textbf{four edges}. Its length is 4. A--C--D--E--G is another four-edge route. Since no shorter path exists, the \textbf{shortest-path distance} from A to G is 4.

A network is \textbf{connected} when every location can reach every other location through some path. Direct contact is not required: A and G have no shared edge, but there is a path between them.
\begin{pointbox}{Distance is a model choice}
In this unweighted network, \emph{distance} counts links, not miles or minutes. We can change the model to include travel time or another number on each connection. That number is an \textbf{edge weight}.
\end{pointbox}
For example, a walk from a station S to a hospital H takes S--C--H (two walkways, 8+9=17 minutes) or S--L--P--H (three walkways, 4+3+5=12 minutes). The route with more edges is faster. A route planner must specify whether it wants the fewest transfers, shortest walking time, or something else.
\begin{trybox}{Try it: read a route}
Why isn't A--D--G a valid path in the county network? Find a path from B to F and count its edges. Does moving the dots without changing the edges alter either answer?
\end{trybox}

\sect{3. A single connection can divide a whole system}
Suppose the connection between the school D and shelter E fails. The communication network separates into two groups: \(\{A,B,C,D\}\) and \(\{E,F,G\}\). Each group is internally reachable, but no path joins the groups. These groups are called \textbf{connected components}. In the original network there is one component; after the failure there are two.
\begin{center}\network{2}\end{center}
A \textbf{bridge} is an edge whose removal increases the number of connected components. D--E is a bridge. B--C is not: if B--C disappears, B can still communicate with C through A or D. The idea is useful whenever a community depends on one road, one partnership, or one communication channel.
\begin{trybox}{Try it: a fragile partnership}
Return to the five organizations with edges PL, PS, LS, SF, FH. Which partnership leaves the senior hub cut off if removed? Is the connection between the school and fridge also a bridge? Trace an alternate path, if one exists.
\end{trybox}

\sect{4. Improve the network with one additional connection}
Suppose funding allows \emph{one} new two-way contact. We compare two proposals:
\begin{center}\begin{minipage}{.48\linewidth}\centering\network{1}\\\small\textbf{A: Add A--D}\\Makes contacting D easier.\end{minipage}\hfill\begin{minipage}{.48\linewidth}\centering\network{3}\\\small\textbf{B: Add D--F}\\Creates a second cross-region connection.\end{minipage}\end{center}
Both proposals increase the number of edges from nine to ten. But the outcomes differ. A--D shortens a local route and leaves D--E as the only link connecting the regions. D--F gives the right-hand group another way to be reached when D--E fails.
\begin{center}\network{4}\end{center}
With D--F added and D--E failed, A can still reach G using A--B--D--F--G. The new connection provides \textbf{redundancy}: another way to carry out a task when one link is lost.

The repair does not solve every problem. If the school D itself stops relaying, both cross-region routes through D disappear. Our conclusion is specific: D--F protects against failure of the \emph{link} D--E, not failure of the \emph{location} D.
\begin{pointbox}{Choosing the best design requires an objective}
If the county values one-step contact from A to D, adding A--D is helpful. If it values keeping the two regions in communication after D--E fails, adding D--F is better. Equal numbers of connections do not imply equal performance. \textbf{The arrangement matters.}
\end{pointbox}

\sect{5. Watch an alert spread through the connections}
Now keep the original nine links. At round 0 only A knows an alert. In every round, each informed location sends it to direct neighbors. A newly informed location can relay it \emph{in the following round}. Every link works, and a location stays informed after it receives the alert.
\begin{center}\small\begin{tabular}{@{}clcl@{}}\toprule Round & Newly informed & Total informed & Total count\\\midrule 0 & A initially & A & 1\\1 & B, C & A, B, C & 3\\2 & D & A, B, C, D & 4\\3 & E & A, B, C, D, E & 5\\4 & F, G & all seven & 7\\\bottomrule\end{tabular}\end{center}
This is a simple example of \textbf{spread} or diffusion. Under this precise rule, the first round when a node receives the alert equals its shortest-path distance from A. Multiple messages arriving at D still inform only one new location. And if D--E is removed, alert rounds cannot reach the other component, no matter how long we wait.
\begin{trybox}{Try it: distinguish counts}
At round 2, how many locations are newly informed, and how many are informed in total? What changes if the new D--F link is added? Under what assumption are these round counts accurate?
\end{trybox}

\sect{6. Transfer the method to another question}
A study group, neighborhood association, or food-sharing network has people and organizations in place of communication sites. Draw each direct contact. Ask a precise question. Then use the same methods to count neighbors, trace paths, test an outage, or compare a proposed connection.

Some behaviors require an additional \textbf{decision rule}. A person may wait until two friends have joined a group before joining; an emergency alert, by contrast, can be relayed as soon as it arrives. The connections tell us \emph{who can influence whom}. The update rule tells us \emph{when something actually changes}. The optional AI application shows how software uses the same kind of node-and-edge representation to plan routes.

\begin{pointbox}{A useful five-step network investigation}
\begin{enumerate}[leftmargin=1.5em,itemsep=1pt]
\item Name what the nodes and edges represent.
\item Identify what you want to find out.
\item Count or trace the relevant connections.
\item Change one feature and compare the result.
\item Explain what the simplified model leaves out.
\end{enumerate}
\end{pointbox}
\sect{Chapter review}
Use the original county network unless a question specifies something different.
\begin{enumerate}
\item List D's direct neighbors and degree. Why does the sum of all seven degrees equal 18 rather than nine?
\item Write one shortest path from A to G and give its length. Explain why A and G are connected without being adjacent.
\item What are the components after D--E fails? Give one edge that is not a bridge and explain how you know.
\item The county can add A--D or D--F. Which proposal keeps A connected to G when D--E fails? Supply an actual path.
\item Under the round-by-round alert rule, list the newly informed locations in rounds 1, 2, 3, and 4. What round first informs E?
\item In the five-partner food network PL, PS, LS, SF, FH, identify one bridge and one node whose loss would isolate H.
\item The hospital route with three walkways takes 12 minutes while the one with two takes 17. How can both facts be true?
\item Describe a real network from your experience. State what a node is, what an edge is, and one limitation of treating every edge equally.
\end{enumerate}
\begin{pointbox}{Answers to the calculation questions}
1. B, C, E; degree 3. Each edge has two ends and contributes twice to the degree total. 2. A--B--D--E--G; length 4. A path exists despite no direct A--G edge. 3. \(\{A,B,C,D\}\) and \(\{E,F,G\}\). For example B--C is not a bridge because B--A--C remains. 4. D--F; A--B--D--F--G. 5. B,C; D; E; F,G. E is first informed in round 3. 6. F--H is a bridge; removal of F isolates H. 7. Each edge can take a different amount of time; number of walkways differs from total minutes. 8. Answers vary.
\end{pointbox}
\end{document}
